\section{Algorithms}
\label{sec:algorithms}

In this section we describe specialized procedures for evaluating the valid possible worlds for selections involving aggregate comparisons.

% ProvSQL implements multiset semantics, using \emph{tuple occurrences} carrying globally unique provenance identifiers (UUIDs).
% We maintain a uniqueness invariant for every intermediate query result $\angsem{\hat I}{q}$: no pair $(u,\alpha)$ occurs twice, meaning that duplicates of the same value-tuple $u$ are represented as distinct occurrences with different annotations $\alpha$.
% Under this invariant, the multiset of group member occurrences can be treated as a finite \emph{set}, and possible worlds are then ordinary subsets.
% This does not change our semantics: multiset properties are preserved because distinct occurrences may still yield equal values $t(u)$, so the induced aggregate input $\vals_t(W)=\mset{t(u)\mid (u,\alpha)\in W}$ remains a multiset of values.

All our algorithms are applied \emph{per group key} $\vec v$, on the
$\preceq$-ordered sequence of group-member
occurrences~$U\defeq\nobreak U^{\preceq}_{\angsem{\hat I}{q'},\vec v}$ of the
\emph{pre-aggregation} sub-query $q'$
(Definition~\ref{def:possible_world}); the order being fixed, we
identify a possible world $W\sqsubseteq U$ with its set of positions, an
index set $W\subseteq [|U|]$. Terms are evaluated on tuples produced
by~$q'$, whose arity is denoted~$k$.

We first provide a dynamic programming algorithm for \sql{SUM} comparisons that incrementally constructs all possible worlds producing the desired sum, named \textsc{ValidSum}.
It processes each tuple iteratively and maintains, for every possible value of the sum, the collection of worlds that realize it.
For a new tuple, previously constructed possible worlds can either remain unchanged or be extended by inclusion of this new tuple, hence creating new worlds having larger sums.
After all the tuples have been encountered, the algorithm returns a collection of non-empty worlds whose total sum satisfies the desired comparison.
The full listing of this algorithm is given in the appendix.

The algorithm works also for the special case where $t=1$, i.e., $\text{\sql{COUNT(*)}}\op C$ comparisons.
However, for this case we can use a more efficient enumeration algorithm, named \textsc{ValidCount}, exploiting the unit-weight sum structure of \sql{COUNT} aggregates.
For $\text{\sql{COUNT(*)}}\op C$ we only rely on the fact that the result of a \sql{COUNT} depends only on the number of occurrences of the group key values of the selected tuples.
Our algorithm recursively constructs worlds by deciding for each tuple-occurrence if it is included in or excluded from the current world.
This enumerates subsets of the desired sizes.
Finally, depending on the comparison operator, we combine the appropriate size-ranges to get a collection of non-empty worlds that satisfy the comparison.
The full listing of this algorithm is given in the appendix.

For other aggregate functions, using brute-force enumeration of non-empty worlds $W\sqsubseteq\nobreak U$ will always give the correct answer by our semantics.
When a side of the comparison mentions several aggregates of the group, as in \sql{SUM(a) > 2 * COUNT(*)}, the valid worlds are those $W$ on which the two sides, both read in~$W$, compare as required; \textsc{ValidSum} and \textsc{ValidCount} do not apply, and our implementation enumerates the worlds by brute force.
\begin{toappendix}
\label{appendix:algorithms}
This appendix gives the full listings of \textsc{ValidSum}
(Algorithm~\ref{alg:sum_dp}) and \textsc{ValidCount}
(Algorithm~\ref{alg:count_enum}), then proves the results of
Section~\ref{sec:algorithms}.
\begin{algorithm}
  \small
  \caption{Enumeration of valid worlds for \sql{SUM}$(t)\op C$ at a fixed group key $\vec v$}
  \label{alg:sum_dp}

  \nlset{}\textbf{\textsc{ValidSum}}:\;

  \Input{a finite set of group-member occurrences $U \defeq U^{\preceq}_{\angsem{\hat I}{q'},\vec v}$, an arbitrary ordering $\preceq$ of $U$, a term $t$ of max-index $\le k$ such that $t(u)\in\NN$ for every $(u,\alpha)\in U$, a constant $C\in\NN$, and an operator ${\op}\in\{=,>,<,\ge,\le,\neq\}$}
  \Output{$\mathcal W =\{\,W\sqsubseteq U\mid W\neq\emptyset,\mathrm{agg}_{t,f_\Sigma}(W)\op C\,\}$}
  \BlankLine
  \Fn{\textsc{ValidSum}$(U,t,C,\op, \preceq)$}{
  \If{$\op$ is $\neq$}{
    \Return $\textsc{ValidSum}(U,t,C,<, \preceq) \cup \textsc{ValidSum}(U,t,C,>, \preceq)$\;
  }
 \BlankLine
  Let $((u_r,\alpha_r))_{r \in [|U|]}$ be the occurrences of $U$ ordered by $\preceq$\\
  $T \gets \mathrm{agg}_{t,f_\Sigma}([|U|]) = \sum_{r=1}^{|U|} t(u_r)$\;
  \If{($\op$ is $=$ \textbf{and} $C > T$) \textbf{or} ($\op$ is $>$ \textbf{and} $C \ge T$) \textbf{or} ($\op$ is $\ge$ \textbf{and} $C > T$) \textbf{or} ($\op$ is $<$ \textbf{and} $C = 0$)}{
    \Return $\emptyset$\;
  }
  \If{${\op}\in\{>,\ge\}$}{
    $J \gets T$\;
  }
  \ElseIf{$\op$ is $<$}{
    $J \gets \min(C-1,\ T)$\;
  }
  \Else{
    $J \gets \min(C,\ T)$\;
  }

  \BlankLine
  Initialize $dp[0] \leftarrow \{\emptyset\}$\;
  \For{$j \gets 1$ \KwTo $J$}{
    $dp[j] \leftarrow \emptyset$\;
  }

  \BlankLine
  $\mathrm{pref\_sum} \gets 0$\;

  \For{$i \gets 1$ \KwTo $|U|$}{
    $w \gets t(u_i)$\;
    $\mathrm{pref\_sum} \gets \mathrm{pref\_sum} + w$\;
    $j_{\max} \gets \min(J,\ \mathrm{pref\_sum})$\;

    \For{$j \gets j_{\max}, j_{\max}-1, \dots, w$}{
      $p \gets j - w$\;
      \If{$dp[p] \neq \emptyset$}{
        let $s \gets |dp[p]|$\;
        \For{$r \gets 1$ \KwTo $s$}{
          append $dp[p][r] \cup \{i\}$ to $dp[j]$\;
        }
      }
    }
  }

  \BlankLine
  \If{$\op$ is $=$}{
    $\mathcal W \leftarrow dp[C]$\;
  }
  \ElseIf{$\op$ is $>$}{
    $\mathcal W \leftarrow \bigcup_{j=C+1}^{J} dp[j]$\;
  }
  \ElseIf{$\op$ is $<$}{
    $\mathcal W \leftarrow \bigcup_{j=0}^{C-1} dp[j]$\;
  }
  \ElseIf{$\op$ is $\ge$}{
    $\mathcal W \leftarrow \bigcup_{j=C}^{J} dp[j]$\;
  }
  \Else{
    $\mathcal W \leftarrow \bigcup_{j=0}^{C} dp[j]$\;
  }

  $\mathcal W \leftarrow \mathcal W \setminus \{\emptyset\}$\;

  \Return $\mathcal W$\;
  }
\end{algorithm}
\end{toappendix}


The following theorem states that \textsc{ValidSum} and
\textsc{ValidCount} return exactly the index set of the $\oplus$-sum of
Definition~\ref{def:semantics}; the predicate provenance is then the
$\oplus$-sum of $\ann_U(W)$ over the returned worlds.
\begin{toappendix}
\begin{algorithm}
  \small
  \caption{Enumeration of valid worlds for \sql{COUNT}$(t)\op C$ at a fixed group key $\vec v$}
  \label{alg:count_enum}

  \nlset{}\textbf{\textsc{ValidCount}:}\;

  \Input{a finite set of group-member occurrences $U\defeq U^{\preceq}_{\angsem{\hat I}{q'},\vec v}$, a constant $C\in\NN$, and an operator ${\op}\in \{=,\neq,>,<,\ge,\le\} $}
  \Output{$\mathcal W =\{\,W\sqsubseteq U\mid W\neq\emptyset,|W|\op C\,\}$}
    \BlankLine
  $\mathcal W \gets \emptyset$\;

  \SetKwFunction{Combinations}{Combinations}
  \Fn{\Combinations{$i, x, W$}}{
    \If{$x = 0$}{
      append $W$ to $\mathcal W$\;
      \Return\;
    }
    \If{$i > |U|$}{
      \Return\;
    }
    \If{$|U|-i+1 < x$}{
      \Return\;
    }

    \tcp{Case 1: do not select occurrence $i$}
    \Combinations{$i+1, x, W$}\;

    \tcp{Case 2: select occurrence $i$}
    \Combinations{$i+1, x-1, W \cup \{i\}$}\;
  }

  \SetKwFunction{AddExact}{AddExact}
  \Fn{\AddExact{$x$}}{
    \If{$x\le 0$ \textbf{or} $x>|U|$}{\Return}
    \Combinations{$1, x, \emptyset$}\;
     }

  \BlankLine
  \Switch{$\op$}{
    \Case{$=$}{
      \AddExact{$C$}\;
    }
    \Case{$>$}{
      \For{$x\gets C+1$ \KwTo $|U|$}{\AddExact{$x$}}
    }
    \Case{$\ge$}{
      \For{$x\gets C$ \KwTo $|U|$}{\AddExact{$x$}}
    }
    \Case{$<$}{
      \For{$x\gets 1$ \KwTo $\min(C-1,|U|)$}{\AddExact{$x$}}
    }
    \Case{$\le$}{
      \For{$x\gets 1$ \KwTo $\min(C,|U|)$}{\AddExact{$x$}}
    }
    \Case{$\neq$}{
      \For{$x\gets 1$ \KwTo $|U|$}{
        \If{$x\neq C$}{\AddExact{$x$}}
      }
    }
  }

  \Return $\mathcal W$\;
\end{algorithm}
\end{toappendix}
\begin{theoremrep}[Correctness of the algorithms]\label{thm:algorithms-correct}
\lean{Algorithms/CountEnum}{CountEnum.countEnum_correct}
\lean{Algorithms/SumDP}{SumDP.sumDP_correct}
Let $q'$ be an aggregate-free query of arity~$k$, $\hat I$ a
$\mathbb K$-instance, $\vec v$ a group key, $\preceq$ an ordering,
$U\defeq\nobreak U^{\preceq}_{\angsem{\hat I}{q'},\vec v}$, $C\in\NN$ and
${\op}\in\{=,>,<,\ge,\le,\neq\}$. Then, on input $(U,C,\op)$ and
$(U,t,C,\op,\preceq)$ respectively, \textsc{ValidCount} and
\textsc{ValidSum} output the set
$\{\,W\sqsubseteq U\mid W\neq\emptyset,\ \mathrm{agg}_{t,f_\Sigma}(W)\op C\,\}$
of valid worlds, where $t$ is the constant term~$1$ for
\textsc{ValidCount} and, for \textsc{ValidSum}, any term of max-index
$\le k$ taking values in $\NN$ on the occurrences of~$U$; $f_\Sigma$
is the \sql{SUM} aggregate, so that $\mathrm{agg}_{1,f_\Sigma}(W)=|W|$.
\end{theoremrep}

\begin{appendixproof}
\paragraph*{Correctness of \textsc{ValidSum}}
Write
\[
  T\defeq \mathrm{agg}_{t,f_\Sigma}([|U|])=\sum_{r=1}^{|U|} t(u_r).
\]

\inlinepar{Case $\op$ is $\neq$.}
If $\op$ is $\neq$, the algorithm returns
\[
\textsc{ValidSum}(U,t,C,<, \preceq)\cup\textsc{ValidSum}(U,t,C,>, \preceq).
\]
By the correctness proved below for ${\op}\in\{<,>\}$, this equals
\[
  \{\,W\subseteq [|U|]\mid W\neq\emptyset,\sum_{r\in W}t(u_r)<C\,\}
\cup
\{\,W\subseteq [|U|]\mid W\neq\emptyset,\sum_{r\in W}t(u_r)>C\,\}.
\]
Since for every integer $x$ one has $x\neq C$ iff $x<C$ or $x>C$, the latter union is exactly
\[
  \{\,W\subseteq [|U|]\mid W\neq\emptyset,\sum_{r\in W}t(u_r)\neq C\,\}.
\]
Hence the theorem holds when $\op$ is $\neq$.

\smallskip\noindent
The remainder of the proof establishes correctness for ${\op}\in\{=,>,<,\ge,\le\}$.

\inlinepar{Early-return cases.}
The algorithm returns $\emptyset$ in the following scenarios:
\begin{itemize}
\item if $\op$ is $=$ and $C>T$;
\item if $\op$ is $>$ and $C\ge T$;
\item if $\op$ is $\ge$ and $C>T$;
\item if $\op$ is $<$ and $C=0$.
\end{itemize}
If $\op$ is $=$ and $C>T$, then every world $W\subseteq [|U|]$ satisfies
\[
  \sum_{r\in W} t(u_r)\le \sum_{r=1}^{|U|} t(u_r)=T<C,
\]
so no world has sum $C$.

If $\op$ is $>$ and $C\ge T$, then every world $W\subseteq [|U|]$ satisfies
\[
  \sum_{r\in W} t(u_r)\le T\le C,
\]
so no world has sum strictly greater than $C$.

If $\op$ is $\ge$ and $C>T$, then again every world has sum at most $T<C$, so no world has sum at least $C$.

If $\op$ is $<$ and $C=0$, since all values $t(u_r)$ lie in $\NN$, every world has nonnegative sum; hence no world has sum $<0$.

Therefore, in each early-return case, the algorithm correctly outputs $\emptyset$.

\inlinepar{DP invariant.}
We can safely assume from now on that none of the early-return cases occur.
For $0\le i\le |U|$, $j\ge 0$, define
\[
\mathcal D_i(j)\defeq
\{\,W\subseteq [i]\mid \sum_{r\in W} t(u_r)=j\,\}.
\]
We claim that after completing the outer-loop iteration with index $i$,
for every $j\in\{0,\dots,J\}$,
\begin{equation}\label{eq:sumdp-invariant}
dp[j]=\mathcal D_i(j).
\end{equation}

We prove this by induction on $i$.

\smallskip\noindent\emph{Base case $i=0$.}
The algorithm initializes
\[
dp[0]=\{\emptyset\}
\qquad\text{and}\qquad
dp[j]=\emptyset \text{ for } 1\le j\le J.
\]
We see,
\[
\mathcal D_0(0)=\{\emptyset\}
\qquad\text{and}\qquad
\mathcal D_0(j)=\emptyset \text{ for } j\ge 1.
\]
Hence \eqref{eq:sumdp-invariant} holds for $i=0$.

\smallskip\noindent\emph{Inductive step.}
Assume \eqref{eq:sumdp-invariant} holds after iteration $i-1$, for some $1\le i\le |U|$.
Fix $j\in\{0,\dots,J\}$.
We show that after iteration $i$, one has $dp[j]=\mathcal D_i(j)$.

We now establish two cases.

\smallskip\noindent\underline{Case 1: $j<t(u_i)$.}
The inner loop does not update $dp[j]$, since the loop runs only for indices
\[
  j=j_{\max},j_{\max}-1,\dots,t(u_i).
\]
Hence after iteration $i$, the value of $dp[j]$ is unchanged, so by the induction hypothesis
\[
dp[j]=\mathcal D_{i-1}(j).
\]
Further, no subset of $[i]$ that contains $i$ can have a sum $j$, because such a subset has a sum at least
\[
  t(u_i)>j,
\]
all summands being nonnegative. Therefore every world in $\mathcal D_i(j)$ must avoid $i$, so
\[
\mathcal D_i(j)=\mathcal D_{i-1}(j).
\]
Consequently,
\[
dp[j]=\mathcal D_i(j).
\]

\smallskip\noindent\underline{Case 2: $j\ge t(u_i)$.}
A subset $W\subseteq [i]$ satisfies
\[
  \sum_{r\in W} t(u_r)=j
\]
iff exactly one of the following two disjoint possibilities hold:
\begin{itemize}
  \item $i\notin W$, in which case $W\subseteq [i-1]$ and
\[
W\in \mathcal D_{i-1}(j);
\]
\item $i\in W$, in which case
\[
W=W'\cup\{i\}
\]
for a unique subset $W'\subseteq [i-1]$ satisfying
\[
  \sum_{r\in W'} t(u_r)=j-t(u_i),
\]
that is,
\[
  W'\in \mathcal D_{i-1}(j-t(u_i)).
\]
\end{itemize}

By the induction hypothesis, before processing index $i$ we have
\[
dp[j]=\mathcal D_{i-1}(j)
\qquad\text{and}\qquad
dp[j-t(u_i)]=\mathcal D_{i-1}(j-t(u_i)).
\]
During the update for this fixed $j$, the algorithm sets
\[
  p\gets j-t(u_i),
\qquad
s\gets |dp[p]|.
\]
It then appends
\[
dp[p][r]\cup\{i\}
\qquad (1\le r\le s)
\]
to $dp[j]$.

Thus the worlds added to $dp[j]$ are exactly those obtained by adding $i$ to the worlds that were already present in $dp[p]$ before any append operation at this $j$.
Since, by the I.H, those pre-existing worlds are exactly the elements of $\mathcal D_{i-1}(j-t(u_i))$, the algorithm adds exactly the worlds in
\[
  \{\,W'\cup\{i\}\mid W'\in \mathcal D_{i-1}(j-t(u_i))\,\}.
\]

$s=|dp[p]|$ is essential when $t(u_i)=0$, for then $p=j$.
In that situation the algorithm appends new worlds to $dp[j]$ while scanning from $dp[j]$, but it scans only the first $s$ entries, i.e., exactly those worlds that were present before the append operations started.
Hence each world in $\mathcal D_{i-1}(j)$ contributes exactly one extended world containing $i$, and none of the newly appended worlds is reused again during the same update.

Therefore, after iteration $i$, the content of $dp[j]$ is exactly:
\begin{itemize}
\item all worlds in $\mathcal D_{i-1}(j)$, corresponding to the case $i\notin W$;
\item together with all worlds of the form $W'\cup\{i\}$ for
  $W'\in\mathcal D_{i-1}(j-t(u_i))$, corresponding to the case $i\in W$.
\end{itemize}
By the disjoint case analysis above, this is precisely $\mathcal D_i(j)$.
Thus
\[
dp[j]=\mathcal D_i(j).
\]

Since both cases have been established, \eqref{eq:sumdp-invariant} holds after iteration $i$.
This completes the proof by induction.

Hence, after the outer loop terminates, for every $j\in\{0,\dots,J\}$,
\[
dp[j]
=
\{\,W\subseteq [|U|]\mid \sum_{r\in W}t(u_r)=j\,\}.
\]

\inlinepar{Selecting satisfying worlds.}
The algorithm forms $\mathcal W$ as a union of those $dp[j]$ whose index $j$ satisfies the comparison with $C$:
\begin{itemize}
\item if $\op$ is $=$, it takes $dp[C]$;
\item if $\op$ is $>$, it takes $\bigcup_{j=C+1}^{J} dp[j]$;
\item if $\op$ is $<$, it takes $\bigcup_{j=0}^{C-1} dp[j]$;
\item if $\op$ is $\ge$, it takes $\bigcup_{j=C}^{J} dp[j]$;
\item if $\op$ is $\le$, it takes $\bigcup_{j=0}^{C} dp[j]$.
\end{itemize}
By the characterization of $dp[j]$ proved above, each such union is exactly the set of worlds
$W\subseteq [|U|]$ whose sum satisfies
\[
  \sum_{r\in W} t(u_r)\op C.
\]
The only possible extra world is $\emptyset$, which belongs to $dp[0]$ and may therefore appear whenever the selected range contains $0$.
The algorithm omits it by setting
\[
\mathcal W\gets \mathcal W\setminus \{\emptyset\}.
\]
It follows that the final output is exactly
\[
  \{\,W\subseteq [|U|]\mid W\neq\emptyset,\sum_{r\in W} t(u_r)\op C\,\}.
\]

This establishes the correctness.

\paragraph*{Correctness of \textsc{ValidCount}}
The computation for \sql{COUNT} depends only on the occurrences in $U$.

\inlinepar{Correctness of \textsc{Combinations}.}
For integers $i$ and $x$, and for a $W\subseteq [|U|]$, consider a call
\[
\textsc{Combinations}(i,x,W).
\]
We claim that this call appends to $\mathcal W$ exactly the sets
\[
  W'\subseteq [|U|]
\]
such that:
\begin{itemize}
\item $W\subseteq W'$;
\item every element of $W'\setminus W$ belongs to $\{i,\dots,|U|\}$;
\item $|W'\setminus W|=x$.
\end{itemize}
Equivalently, it appends exactly the worlds of the form
\[
W\cup S
\quad\text{with}\quad
S\subseteq \{i,\dots,|U|\}
\text{ and } |S|=x.
\]

We prove this by induction on the pair $(|U|-i+1,x)$.

\smallskip\noindent\emph{Base cases.}
If $x=0$, it appends $W$ and returns. This is correct because the only way to extend $W$ by choosing exactly $0$ further occurrences is to keep $W$.

If $i>|U|$, there are no remaining occurrences to choose from. Hence no extension of $W$ by a positive number of further occurrences exists, so returning without appending anything is correct.

If $|U|-i+1<x$, then fewer than $x$ occurrences remain available. Thus it is impossible to choose exactly $x$ more occurrences, so returning without appending is again correct.

\smallskip\noindent\emph{Inductive step.}
Assume now that $x>0$, $i\le |U|$, and $|U|-i+1\ge x$.
Any extension of $W$ by exactly $x$ occurrences chosen from
\[
\{i,\dots,|U|\}
\]
falls into one of two disjoint cases:
\begin{itemize}
\item it does \emph{not} contain $i$;
\item it \emph{does} contain $i$.
\end{itemize}

In the first case, the required worlds are exactly those produced by the recursive call
\[
\textsc{Combinations}(i+1,x,W),
\]
since one must still choose all $x$ required occurrences from the suffix starting at $i+1$.

In the second case, the desired worlds are exactly those produced by
\[
\textsc{Combinations}(i+1,x-1,W\cup\{i\}),
\]
since after selecting $i$, one must select exactly $x-1$ further occurrences from the suffix starting at $i+1$.

By the I.H, the first recursive call appends exactly the worlds of the first kind, and the second appends exactly the worlds of the second kind. Since the two cases are disjoint and exhaustive, \textsc{Combinations} appends exactly all desired worlds, each exactly once.

This proves the claim.

\inlinepar{Correctness of \textsc{AddExact}.}
Consider a call \textsc{AddExact}$(x)$.

If $x\le 0$, the procedure returns immediately. This is correct because the output requires $W\neq\emptyset$, so no world of size $0$ should be output, and no world has negative cardinality.

If $x>|U|$, the procedure also returns immediately. This is correct because no subset of $[|U|]$ can have cardinality greater than $|U|$.

If $1\le x\le |U|$, the procedure calls
\[
\textsc{Combinations}(1,x,\emptyset).
\]
By the claim proved above, this appends exactly all sets
\[
  W\subseteq [|U|]
\quad\text{such that}\quad
|W|=x.
\]
Hence \textsc{AddExact}$(x)$ appends exactly the non-empty worlds of cardinality $x$.

\inlinepar{Selecting satisfying worlds.}
We now inspect the final case distinction on $\op$.

\begin{itemize}
  \item If $\op$ is $=$, the algorithm executes \textsc{AddExact}$(C)$, so by the correctness of \textsc{AddExact} it outputs exactly the non-empty worlds $W\subseteq [|U|]$ such that $|W|=C$.

\item If $\op$ is $>$, the algorithm executes \textsc{AddExact}$(x)$ for every
\[
x\in\{C+1,\dots,|U|\}.
\]
Therefore it returns exactly the non-empty worlds whose cardinality is strictly greater than $C$.

\item If $\op$ is $\ge$, the algorithm executes \textsc{AddExact}$(x)$ for every
\[
x\in\{C,\dots,|U|\}.
\]
Since \textsc{AddExact} ignores values $x\le 0$, this outputs exactly the non-empty worlds whose cardinality is at least $C$.

\item If $\op$ is $<$, the algorithm executes \textsc{AddExact}$(x)$ for every
\[
x\in\{1,\dots,\min(C-1,|U|)\}.
\]
Thus it returns exactly the non-empty worlds whose cardinality is strictly less than $C$.

\item If $\op$ is $\le$, the algorithm executes \textsc{AddExact}$(x)$ for every
\[
x\in\{1,\dots,\min(C,|U|)\}.
\]
Thus it returns exactly the non-empty worlds whose cardinality is at most $C$.

\item If $\op$ is $\neq$, the algorithm executes \textsc{AddExact}$(x)$ for every
\[
x\in\{1,\dots,|U|\}\setminus\{C\}.
\]
Hence it returns exactly the non-empty worlds whose cardinality is different from $C$.
\end{itemize}

In every case, the resulting  $\mathcal W$ is exactly
\[
  \{\,W\subseteq [|U|]\mid W\neq\emptyset,\ |W|\op C\,\}.
\]

This establishes correctness.
\end{appendixproof}


We identify certain aggregate-comparison predicates for which the valid worlds can be constructed in polynomial time:
\begin{propositionrep}
  \label{prop:algorithms}
  The provenance of \sql{HAVING} queries of the following forms can be computed
  in polynomial time in data complexity:
  \begin{itemize}
    \item $\text{\sql{MIN(a)}}\op c$,
    $\text{\sql{MAX(a)}}\op c$,
$\text{\fakesql{PICKFIRST}\texttt{(a)}}\op c$
      with
      ${\op}\in\{<,\leq,=,\geq,>,\neq\}$ in absorptive m-semirings in which
      $\otimes$ distributes over $\ominus$;
    \item \sql{COUNT(*) >= k} and \sql{COUNT(*) > k} for $k$ a fixed constant, in absorptive
      m-semirings;
\item \sql{SUM(a) >= }$c$ and \sql{SUM(a) > }$c$ with $c\in\NN$, for an $\NN$-valued
  attribute \texttt{a} in which every nonzero value $a$ satisfies
  $\frac ca\leq k$ where $k\geq 0$ is a fixed constant, in absorptive
      m-semirings;
    \item \sql{COUNT(*) = k}, \sql{COUNT(*) <= k}, and \sql{COUNT(*) < k} for $k$ a fixed
      constant, in arbitrary m-semirings;
    \item \sql{SUM(a) = }$c$, \sql{SUM(a) <= }$c$, and \sql{SUM(a) < }$c$
      with $c\in\NN$, for an $\NN$-valued attribute \texttt{a} whose
      values are all nonzero and satisfy $\frac ca\leq k$ where $k\geq 0$
      is a fixed constant, in arbitrary m-semirings.
  \end{itemize}
\end{propositionrep}
\begin{proof}
Fix an arbitrary $\mathbb K$-instance $\hat I$ and a group key $\vec v$, and write $U\defeq U^{\preceq}_{\angsem{\hat I}{q'},\vec v}$ for the sequence of group-member occurrences at~$\vec v$.
For each case below, we show that the provenance for $\vec v$ can be computed in time polynomial in $|U|$. Since the queries are fixed, this implies polynomial time in data complexity.
We recall that the provenance of the \sql{HAVING} condition (equiv.\ selection predicates with aggregate comparisons) at group key $\vec v$ is obtained by enumerating the valid non-empty worlds
and then forming the corresponding semiring sum of the $\mathbb K$-annotations associated with these worlds using the possible-world semantics. Thus it suffices to show that, in each of the stated cases, the relevant list of worlds can be computed in polynomial time.

\inlinepar{Case 1: $\text{\sql{MIN(a)}}\op c$, $\text{\sql{MAX(a)}}\op c$, and $\text{\fakesql{PICKFIRST}\texttt{(a)}}\op c$.}
Write $\alpha_i$ and $a_i$ for the annotation and the value of
attribute~\texttt{a} of the $i$-th occurrence of~$U$, occurrences being
indexed in the order~$\preceq$, and
$L\defeq\bigoplus_{a_i<c}\alpha_i$, $L'\defeq\bigoplus_{a_i\le c}\alpha_i$,
$G\defeq\bigoplus_{a_i\ge c}\alpha_i$, $G'\defeq\bigoplus_{a_i>c}\alpha_i$,
$E\defeq\bigoplus_{a_i=c}\alpha_i$. In an absorptive m-semiring in which
$\otimes$ distributes over $\ominus$, the possible-world sum for
$\text{\sql{MIN(a)}}\op c$ has the closed form
\lean{HavingMinMax}{Having.minScan_correct}
\[
\begin{array}{r@{\ }l@{\qquad}r@{\ }l@{\qquad}r@{\ }l}
{<}c: & L, & {\le}c: & L', & {\ne}c: & L\oplus(\mathbb 1\ominus L')\otimes G',\\
{\ge}c: & (\mathbb 1\ominus L)\otimes G, & {>}c: & (\mathbb 1\ominus L')\otimes G', & {=}c: & (\mathbb 1\ominus L)\otimes E;
\end{array}
\]
for instance, $\text{\sql{MIN(a)}}<c$ holds in a world exactly when some
occurrence with $a_i<c$ is present, and absorptivity reduces the sum over
these worlds to the sum over the singletons, whereas $\text{\sql{MIN(a)}}\ge c$
additionally requires that no occurrence with $a_i<c$ be present, whence
the factor $\mathbb 1\ominus L$. The closed form for $\text{\sql{MAX(a)}}\op c$
is the mirror image, exchanging $<$ and~$>$
\lean{HavingMinMax}{Having.maxScan_correct}, and for
$\text{\fakesql{PICKFIRST}\texttt{(a)}}\op c$ (the non-commutative aggregate of
Example~\ref{ex:aggregates}) it is
$\bigoplus_{i:\,a_i\op c}(\mathbb 1\ominus\bigoplus_{j<i}\alpha_j)\otimes\alpha_i$,
one term per satisfying occurrence that comes first
\lean{HavingMinMax}{Having.firstScan_correct}. Each is computed by a
single scan of~$U$ using $O(|U|)$ semiring operations. Both hypotheses
are used: absorptivity fails in the tropical m-semiring over
$\mathbb R\cup\{\infty\}$ (idempotent and distributive), where the scan
value differs from the possible-world provenance already for
$\text{\sql{MIN(a)}}\ge c$ on a two-occurrence group
\lean{Semirings/Tropical}{TropicalR.minScan_ne_prov}; distributivity of
$\otimes$ over $\ominus$ identifies the world annotation $\ann_U(W)$ with
the expanded form $T_U(W)$ of Section~\ref{sec:results}
\lean{HavingSemantics}{Having.worldAnn_eq_T}, on which the closed forms
are established.

\inlinepar{Case 2: \sql{COUNT(*) >= k} and \sql{COUNT(*) > k} for fixed $k$, in absorptive m-semirings.}
By correctness of Algorithm~\ref{alg:count_enum}
(Theorem~\ref{thm:algorithms-correct}), the valid worlds are exactly the non-empty worlds
$W\sqsubseteq U$ such that $|W|\ge k\text{ or }|W|>k$, respectively.
  If such a $W$ satisfies the condition, then every superset of $W$ also satisfies it; equivalently these $W$ form an upward-closed family of sets. The m-semiring $\mathbb K$ being absorptive, the possible-world provenance of an upward-closed family of worlds (the $\oplus$-sum of their world annotations) is equal to the $\oplus$-sum of the monomials of its minimal worlds \lean{Having}{Having.upward_closed_collapse}; for the family of worlds of size at least $k$ this collapse is exactly the identity $F_k(U)=S_k(U)$ of Section~\ref{sec:results} \lean{Having}{Having.F_eq_S}. Here, the minimal satisfying worlds are:
\begin{itemize}
\item exactly the worlds of size $k$ for \sql{COUNT(*) >= k};
\item exactly the worlds of size $k+1$ for \sql{COUNT(*) > k}.
\end{itemize}
Since $k$ is fixed, the number of such worlds is $\binom{|U|}{k}\text{ or }\binom{|U|}{k+1}$,
which is polynomial in $|U|$. Therefore the provenance can be computed in polynomial time.

\inlinepar{Case 3: \sql{SUM(a) >= c} and \sql{SUM(a) > c} with $c\ge 0$, in absorptive m-semirings.}
By correctness of Algorithm~\ref{alg:sum_dp}
(Theorem~\ref{thm:algorithms-correct}), the valid worlds are exactly the non-empty worlds
$W\sqsubseteq U$ such that $\sum_{r\in W} t(u_r)\ge c$ (resp.\ $>c$).
Since all values are nonnegative, these families are upward-closed under inclusion. Hence, in an absorptive m-semiring, the provenance collapses to the $\oplus$-sum of the monomials of the minimal satisfying worlds
\lean{Having}{Having.sum_ge_collapse} \lean{Having}{Having.sum_gt_collapse}.
Let $W$ be a minimal satisfying world. A zero-valued occurrence never belongs to $W$: removing it leaves the sum unchanged, contradicting minimality. Hence every value $a$ of an occurrence of $W$ is nonzero and satisfies $\frac ca\le k$, i.e., $c\le k\cdot a$, so any $k$ selected occurrences of $W$ already contribute a sum at least $c$. For \sql{SUM(a) >= c}, this gives $|W|\le k$: if $|W|>k$, then any $k$-element subset of $W$ already has sum $\ge c$, contradicting minimality \lean{Having}{Having.minimal_card_le_of_sum_ge}. For the strict comparison \sql{SUM(a) > c}, a $k$-element subset only guarantees a sum $\ge c$, and one further (nonzero) occurrence makes the inequality strict -- covering in particular the boundary case $c=0$ -- so minimality gives $|W|\le k+1$ \lean{Having}{Having.minimal_card_le_of_sum_gt}.

Therefore the number of minimal satisfying worlds is at most
$\sum_{i\le k}\binom{|U|}{i}$ in the first case
\lean{Having}{Having.card_minimal_sum_ge_le} and
$\sum_{i\le k+1}\binom{|U|}{i}$ in the second; both bounds are polynomial
in $|U|$ for fixed~$k$, so the provenance can be computed in polynomial
time.

\inlinepar{Case 4: \sql{COUNT(*) = k}, \sql{COUNT(*) <= k}, and \sql{COUNT(*) < k} for fixed $k$, in arbitrary m-semirings.}
By correctness of Algorithm~\ref{alg:count_enum}, the valid worlds are exactly:
\begin{itemize}
\item the worlds of size exactly $k$ for \sql{COUNT(*) = k};
\item the worlds of sizes $1,\dots,k$ for \sql{COUNT(*) <= k};
\item the worlds of sizes $1,\dots,k-1$ for \sql{COUNT(*) < k}.
\end{itemize}
Since $k$ is fixed, the number of such worlds is bounded by
  $\binom{|U|}{k},\, \sum_{i=1}^{k}\binom{|U|}{i},\text{ and } \sum_{i=1}^{k-1}\binom{|U|}{i}$, respectively \lean{Having}{Having.card_powerset_filter_card_le}. Each of these bounds is polynomial in $|U|$. Therefore the provenance can be computed in polynomial time in arbitrary m-semirings.

\inlinepar{Case 5: \sql{SUM(a) = }$c$, \sql{SUM(a) <= }$c$, and \sql{SUM(a) < }$c$ with all values nonzero and $\frac ca\le k$, in arbitrary m-semirings.}
By correctness of Algorithm~\ref{alg:sum_dp}
(Theorem~\ref{thm:algorithms-correct}), the valid worlds are exactly the
non-empty worlds $W\sqsubseteq U$ whose sum is equal to, at most, or less
than~$c$. If $c=0$ there is none, all values being nonzero. Otherwise, a
valid world $W$ satisfies $|W|\cdot\frac ck\le\sum_{r\in W}t(u_r)\le c$,
since every value is at least $\frac ck$, hence $|W|\le k$. The valid
worlds are therefore among the $\sum_{i=1}^{k}\binom{|U|}{i}$ worlds of
size at most~$k$ \lean{Having}{Having.card_powerset_filter_card_le},
polynomially many for fixed~$k$, and \textsc{ValidSum}, whose table is
bounded by~$c$, enumerates exactly them. No assumption on the m-semiring
is needed.

\smallskip\noindent
With these 5 cases, provenance can be computed in polynomial time for all the \sql{HAVING} queries described in the proposition.
\end{proof}
\begin{proofsketch}
Fix the group key $\vec v$ and write $U\defeq U^{\preceq}_{\angsem{\hat I}{q'},\vec v}$. The provenance of the \sql{HAVING} condition is the semiring sum of the annotations of the valid non-empty worlds, so it suffices to enumerate the relevant worlds in time polynomial in $|U|$.
For \sql{MIN(a)}, \sql{MAX(a)} and \fakesql{PICKFIRST}\texttt{(a)} (Case~1), the validity of a world is decided occurrence by occurrence, so in an absorptive m-semiring with $\otimes$ distributive over $\ominus$ the provenance is obtained by a single scan \lean{HavingMinMax}{Having.minScan_correct}.
For $\text{\sql{COUNT(*)}}\ge k$ or ${}>k$ (Case~2) and $\text{\sql{SUM(a)}}\ge c$ or ${}>c$ (Case~3), the valid worlds form an upward-closed family, whose provenance, in an \emph{absorptive} m-semiring, collapses to the $\oplus$-sum of the monomials of its minimal worlds \lean{Having}{Having.upward_closed_collapse}. These minimal worlds have bounded cardinality: $k$, or $k+1$ for the strict variants (using the hypothesis $\frac ca\le k$ on nonzero values for \sql{SUM}); so there are only $O(|U|^{k+1})$ of them.
Finally, for $\text{\sql{COUNT(*)}}=k$, $\le k$ or $<k$ (Case~4), the valid worlds are exactly those of cardinality at most $k$, again polynomially many, so no absorptivity is needed; likewise for $\text{\sql{SUM(a)}}=c$, $\le c$ or $<c$ when all values are at least $\frac ck$ (Case~5).
\end{proofsketch}

\begin{example}\label{ex:prop-simplification}
Consider as in Example~\ref{ex:semantics} a group of three occurrences,
now with values $a(u_1)=3$ and $a(u_2)=a(u_3)=2$, and the condition
$\text{\sql{SUM(a)}}\ge 5$. The valid worlds are $\{1,2\}$, $\{1,3\}$
and $\{1,2,3\}$, with possible-world provenance
$(x_1\land x_2\land\lnot x_3)\lor(x_1\land x_3\land\lnot x_2)\lor(x_1\land
x_2\land x_3)$; the family is upward-closed and $\mathbb B[X]$
absorptive, so this equals the $\oplus$-sum over the minimal valid
worlds $\{1,2\}$ and $\{1,3\}$, namely $x_1\land(x_2\lor x_3)$.
Proposition~\ref{prop:algorithms} rests on this collapse: the minimal
worlds suffice, polynomially many for a fixed threshold.
\end{example}

The tractability established by Proposition~\ref{prop:algorithms} is
\emph{symbolic}: it computes the provenance of a \sql{HAVING} comparison in an
arbitrary (absorptive) m-semiring by enumerating the relevant valid worlds. A
complementary form of tractability arises in the \emph{probabilistic} setting. By
Proposition~\ref{prop:pqe-equiv}, evaluating a \sql{HAVING} query over a
tuple-independent instance amounts to computing the probability of its
$\mathbb{B}[X]$-provenance under the possible-world semantics. When the
contributors of a group are independent, this probability need not be obtained
by enumerating worlds, whose number is in general exponential: as Ré and
Suciu~\cite{re2009trichotomy} showed in the base case of their safe plans,
the distribution of the aggregate over the group can be assembled by a
convolution of marginal distributions, yielding the answer in polynomial
time for several aggregates and \emph{all} comparison operators. In our
terms:

\begin{propositionrep}[after~\cite{re2009trichotomy}]
\label{prop:agg-cmp-poly-prob}
Let $\hat I$ be a $\mathbb{B}[X]$-instance whose variables carry
independent probabilities, $\vec v$ a group key and
$U\defeq U^{\preceq}_{\angsem{\hat I}{q'},\vec v}=((u_i,\alpha_i))_{1\le i\le N}$
its occurrences; assume the contributors \emph{independent}, i.e., the
$\alpha_i$ have pairwise disjoint sets of variables, disjoint from those
of every other annotation in $\angsem{\hat I}{q}$, and write
$p_i\defeq\Pr(\alpha_i)$. For every comparison operator $\op$ and
threshold $C\in\NN$, the probability of the aggregate comparison on the
group under the possible-world semantics is computable in time
$O\bigl(N\cdot\min(C,N-C)\bigr)$ for $\text{\sql{COUNT(*)}}\op C$, in
time $O(N\cdot\min(C,S))$ for $\text{\sql{SUM}}(t)\op C$ with
$t(u_i)\in\NN$ and $S\defeq\sum_{i=1}^N t(u_i)$, and in time $O(N)$ for
$\text{\sql{MIN}}(t)\op C$ and $\text{\sql{MAX}}(t)\op C$ over any
totally ordered value domain.
\end{propositionrep}

\begin{proof}
By the independence hypothesis and the equivalence between
probabilistic query evaluation and possible-world summation
(Proposition~\ref{prop:pqe-equiv}), the probability of the comparison
atom factors through the indicators
$X_i\defeq[(u_i,\alpha_i)\text{ present in the world}]$, which form a
family of \emph{independent} Bernoullis with $\Pr[X_i=1]=p_i$.
We exhibit the three procedures.

\inlinepar{(1) \sql{COUNT(*)}\,$\op\,C$.}
The aggregate value on a world $W$ is $|W|=\sum_{i=1}^N X_i$, a
Poisson-binomial variable $B$.  For each $\op$, the satisfaction set
$\{j\in\{0,\dots,N\}\mid j\op C\}$ is an interval (possibly truncated
by $0$, or with the singleton $\{C\}$ removed in the $\neq$ case), so
the answer reduces to two CDF queries among $\Pr[B\le T]$,
$\Pr[B\ge T]$, $\Pr[B=T]$, $\Pr[B=0]$.

Let $\rho_n(j)\defeq\Pr\bigl[\sum_{i\le n}X_i=j\bigr]$.
The classical convolution identity
\lean{HavingProbability}{HavingProbability.countMass_insert_succ}
\[
\rho_n(j)=(1-p_n)\,\rho_{n-1}(j)+p_n\,\rho_{n-1}(j-1)
\]
yields a 1-D rolling DP that fills the row $\rho_N(\cdot)$ up to any
prescribed index $J\le N$ in time $O(N\cdot J)$.
Computing the upper tail $\Pr[B\ge C]$ directly costs
$O(N\cdot(N-C+1))$; computing the lower tail costs $O(N\cdot C)$.
Using $\Pr[B\ge C]=\Pr[B'\le N-C]$ for $B'=\sum_i(1-X_i)$, one can
always evaluate the tail on whichever side of $C$ is shorter, giving
the announced bound $O\bigl(N\cdot\min(C,N-C)\bigr)$.
The empty-world mass $\Pr[B=0]=\prod_i(1-p_i)$ is obtained as a
by-product (linear time) and subtracted in the branches whose
satisfaction interval contains $0$ (i.e., $\le$, $<$, $\neq$, and $\ge$
with $C\le0$).

\inlinepar{(2) \sql{SUM}$(t)\,\op\,C$ with integer
weights.}
The aggregate value on a world $W$ is
$\Sigma\defeq\sum_{i:X_i=1}t(u_i)$, a weighted Poisson-binomial
variable.  The same convolution generalizes
\lean{HavingProbability}{HavingProbability.sumMass_insert_of_le}:
\[
\sigma_n(s)=(1-p_n)\,\sigma_{n-1}(s)+p_n\,\sigma_{n-1}\bigl(s-t(u_n)\bigr),
\qquad
0\le s\le S,
\]
where $\sigma_n(s)\defeq\Pr\bigl[\sum_{i\le n}t(u_i)X_i=s\bigr]$.
Only the indices $s\le\min(C,S)$ are needed, the mass of $\Sigma\ge C$
being obtained by complementation (this is the saturation at~$C$ of the
semiring $S_{C+1}$ of~\cite{re2009trichotomy}), so filling
$\sigma_N(\cdot)$ on $[0,\min(C,S)]$ takes $O(N\cdot\min(C,S))$ time and
gives every $\Pr[\Sigma\op C]$ by a single summation; the empty-world
adjustment is the same $\prod_i(1-p_i)$ as before.

\inlinepar{(3) \sql{MIN}$(t)\,\op\,C$ and
\sql{MAX}$(t)\,\op\,C$.}
Independence makes the events
``$\max_{i:X_i=1}t(u_i)\le C$''
\lean{HavingProbability}{HavingProbability.funcProb_maxLeOnNonempty}
and ``$\min_{i:X_i=1}t(u_i)\ge C$''
\lean{HavingProbability}{HavingProbability.funcProb_minGeOnNonempty}
\emph{factor over $i$}, since each surviving $X_i=1$ contributes its
threshold check independently:
\begin{align*}
  \Pr[\text{\sql{MAX}}\le C\text{ on a nonempty world}]
 &=\Bigl(\textstyle\prod_{i:t(u_i)>C}(1-p_i)\Bigr)
   \,\bigl(1-\textstyle\prod_{i:t(u_i)\le C}(1-p_i)\bigr),\\
   \Pr[\text{\sql{MIN}}\ge C\text{ on a nonempty world}]
 &=\Bigl(\textstyle\prod_{i:t(u_i)<C}(1-p_i)\Bigr)
   \,\bigl(1-\textstyle\prod_{i:t(u_i)\ge C}(1-p_i)\bigr),
\end{align*}
each computable by a single scan in $O(N)$.
The remaining operators are obtained by complementation
(``$\text{\sql{MAX}}>C$'' is the negation of ``$\text{\sql{MAX}}\le C$'' on the nonempty
worlds, etc.) and the equality cases by intersecting two such tails;
$\neq$ subtracts the equality from the nonempty mass
$1-\prod_i(1-p_i)$.  All these algebraic combinations stay $O(N)$.
\end{proof}

\begin{toappendix}
The following example illustrates the convolution used in the proof of
Proposition~\ref{prop:agg-cmp-poly-prob}.
\begin{example}\label{ex:poisson-binomial}
Suppose the three contributors of Example~\ref{ex:semantics} are
independent with marginals~$p_1=\tfrac12$, $p_2=\tfrac14$ and
$p_3=\tfrac13$, and consider $\text{\sql{COUNT(*)}}\ge 2$. Writing $\rho_n(j)$
for the probability that exactly $j$ of the first $n$ contributors are
present, so that $\rho_0(0)=1$ (and~$\rho_0(j)=0$ for~$j>0$), the
recurrence $\rho_n(j)=(1-p_n)\,\rho_{n-1}(j)+p_n\,\rho_{n-1}(j-1)$ fills,
row by row, the table
\begin{center}
\begin{tabular}{ccccc}
\toprule
 & $j=0$ & $j=1$ & $j=2$ & $j=3$\\
\midrule
$\rho_0$ & $1$ & & & \\
$\rho_1$ & $\tfrac12$ & $\tfrac12$ & & \\
$\rho_2$ & $\tfrac38$ & $\tfrac12$ & $\tfrac18$ & \\
$\rho_3$ & $\tfrac14$ & $\tfrac{11}{24}$ & $\tfrac14$ & $\tfrac1{24}$\\
\bottomrule
\end{tabular}
\end{center}
Hence
$\Pr[\text{\sql{COUNT(*)}}\ge 2]=\rho_3(2)+\rho_3(3)=\tfrac14+\tfrac1{24}=\tfrac{7}{24}$,
the probability of the Boolean provenance ``at least two of
$x_1,x_2,x_3$'' obtained in Example~\ref{ex:semantics}. This
Poisson-binomial CDF is exactly the quantity our implementation evaluates
for probabilistic \sql{COUNT} comparisons
(Section~\ref{sec:implementation_experiments}).
\end{example}
\end{toappendix}
